\slajdid{02-s010}
\begin{frame}[label=02-s010]
\gusttekst
Na osnovu zahteva koji se postavljaju pred digitalni sistem uobičajeno je da se
\begin{enumerate}
\item Eventualno formira funkcionalna tabela digitalnog sistema
\item Eventualno napišu logičke funkcije
\item Prikaže digitalni sistem u obliku „električne“ šeme upotrebom simbola logičkih funkcija
\end{enumerate}
Upotrebljen je pojam „eventualno“ pošto u sintezi mogu biti preskočeni neki koraci. Na primer bez pisanja logičkih funkcija korišćenjem Bulove algebre moguće je iz funkcionalne table prikazati digitalni sistem upotrebom simbola logičkih funkcija. Isto tako možda možemo da preskočimo i formiranje funkcionalne tabele, pa čak i pisanje logičkih funkcija i da direktno crtamo šemu.

Pojam električna šema je apsolutno ispravan pošto svakom simbolu logičkih funkcija odgovara realno logičko kolo sa svojim pravim električnim karakteristikama.

U ovom procesu sinteze se pojavljuje i među korak a to je „minimizacija“ ili prilagođenje logičkih funkcija zahtevima digitalnog sistema; na primer minimalan broj upotrebljenih tranzistora, minimalno kašnjenje, što manji broj kola iste vrste itd\ldots
\end{frame}
